% 第18章 Schwinger 玻色子平均场理论
\chapter{Schwinger 玻色子平均场理论}

本章综述大 $N$ Schwinger 玻色子平均场理论（SBMFT）在 Heisenberg 铁磁体（FM-B）与反铁磁体（AFM-B）于一维及正方格子上的情形。SBMFT（与第 11 章的自旋波理论相反）保持有效哈密顿量显式的 $SU(2)$ 对称\footnote{M. Takahashi 发展了一个类似 SBMFT 但破坏自旋转动对称的修正自旋波理论。}。我们将看到，SBMFT 恢复第 15 章连续途径的主要结果（$\Theta$ 项效应除外）。

下面把 $N$ 保持为独立参数；计算物理 Heisenberg 模型的自旋关联时取 $N \rightarrow 2$。

SBMFT 由把 (16.39) 作用量中的辅助场换成静态均匀鞍点参数给出：

\begin{equation*}
N \mathcal {S} [ Q _ {ij} ( \tau ), \lambda_ {i} ( \tau ) ] \rightarrow N \mathcal {S} _ {0} ( Q, - i \lambda ) = \beta F ^ {\mathrm{MF}} ( Q, \lambda ).
\tag{18.1}
\end{equation*}

$F^{\mathrm{MF}}$ 是平均场自由能，可写为

\begin{equation*}
F ^ {\mathrm{MF}} ( Q, \lambda ) = - \beta^ {- 1} \ln \operatorname{Tr}_{im} \left[ \exp \left( - \beta H ^ {\mathrm{MF}} [ Q, \lambda ] \right) \right] ,
\tag{18.2}
\end{equation*}

$H^{\mathrm{MF}}$ 是 $N$ 味解耦玻色子的平均场哈密顿量。

\section{铁磁体情形}

用 (18.1) 与 (16.36)，无外源流时 FM-B 的平均场哈密顿量为

\begin{align*}
H ^ {\mathrm{MF}} & = \sum_ {im} \lambda \, a _ {im} ^ {\dagger} a _ {im} - Q \sum_{\langle ij \rangle, m} \left( a _ {im} ^ {\dagger} a _ {jm} + a _ {jm} ^ {\dagger} a _ {im} \right)\\
& \quad + N \mathcal {N} \frac {z Q ^ {2}}{2 J} - N \mathcal {N} S \lambda\\
& = \sum_{\mathbf {k}, m} \epsilon_ {\mathbf {k}} \, a _ {\mathbf {k} m} ^ {\dagger} a _ {\mathbf {k} m} + N \mathcal {N} \frac {z Q ^ {2}}{2 J} - N \mathcal {N} S \lambda .
\tag{18.3}
\end{align*}

$\mathcal{N}$ 与 $z$ 分别是格点数与近邻数，$\mathbf{k}$ 是格子动量，而

\begin{align*}
\epsilon_ {\mathbf {k}} & = \lambda - z Q \gamma_ {\mathbf {k}} ,\\
\gamma_ {\mathbf {k}} & = \frac {1}{z} \sum_ {\hat {\eta}} e ^ {i \mathbf {k} \hat {\eta}} .
\tag{18.4}
\end{align*}

$\eta$ 是近邻矢量。

平均场自由能为

\begin{equation*}
F ^ {\mathrm{MF}} = N \sum_ {\mathbf {k}} \ln \left( 1 - e ^ {- \epsilon_{\mathbf {k}} / T} \right) + N \mathcal {N} \frac {z Q ^ {2}}{2 J} - N \mathcal {N} S \lambda .
\tag{18.5}
\end{equation*}

对 $\lambda$ 与 $Q$ 极小化 $F$，得鞍点方程 (17.1) 与 (17.2)：

\begin{equation*}
\frac {1}{\mathcal {N}} \sum_ {\mathbf {k}} n _ {\mathbf {k}} = S ,
\tag{18.6}
\end{equation*}

\begin{equation*}
\frac {1}{\mathcal {N}} \sum_ {\mathbf {k}} n _ {\mathbf {k}} \gamma_ {\mathbf {k}} = Q / J ,
\tag{18.7}
\end{equation*}

其中

\begin{equation*}
n _ {\mathbf {k}} = \frac {1}{e ^ {\epsilon_{\mathbf {k}} / T} - 1}.
\tag{18.8}
\end{equation*}

方程 (18.6) 与 (18.7) 唯一确定平均场参数 $Q(T, S)$ 与 $\lambda(T, S)$。

用平均场变量更具物理意义的参数化是方便的：

\begin{align*}
\epsilon_ {\mathbf {k}} & = z Q \left( 1 - \gamma_ {\mathbf {k}} + \frac {1}{4z} \kappa^ {2} \right) ,\\
\kappa^ {2} & = 4 z \left( \lambda / z Q - 1 \right) .
\tag{18.9}
\end{align*}

$zQ$、$\kappa$ 分别描述带宽与逆关联长度。图~\ref{fig:18.1} 画出一维铁磁体的平均场色散。

从 (18.6) 减去 (18.7)，得 $Q$ 的显式表达式：

\begin{equation*}
Q = J S - \frac {J}{\mathcal {N}} \sum_ {\mathbf {k}} n _ {\mathbf {k}} \left( 1 - \gamma_ {\mathbf {k}} \right).
\tag{18.10}
\end{equation*}

虽然在小宗量处 $n_{\mathbf{k}} \sim (\beta \epsilon_{\mathbf{k}})^{-1}$，右端被积函数在 $k \to 0$ 有界，求和具有 $T$ 的正则幂级数。低温 $T \ll zQ$ 时：

\begin{equation*}
Q = J S + \mathcal {O} ( T / JS ).
\tag{18.11}
\end{equation*}

\begin{figure}[H]
\centering
\includegraphics[width=0.55\textwidth]{fig18_1.jpg}
\caption{一维铁磁体的平均场色散。}
\label{fig:18.1}
\end{figure}

有限温时 $\kappa > 0$，$\epsilon_{k}$ 在 $k = 0$ 有大小 $zQ\kappa^{2}/8$ 的能隙。取 $T \to 0$、$\mathcal{N} \to \infty$（热力学极限）时，Bose 函数对所有 $k \neq 0$ 消失。为满足约束方程 (18.6)，$k = 0$ 处必须有宏观占据：

\begin{equation*}
S = \frac {1}{\mathcal {N}} \sum_ {\mathbf {k}} n _ {\mathbf {k}} \sim \frac {1}{\mathcal {N}} n _ {0} = \frac {4 T}{\mathcal {N} J S \kappa^ {2}} + \mathcal {O} ( T ^ {2} ) ,
\tag{18.12}
\end{equation*}

这意味着

\begin{equation*}
\lim_{T \to 0} \kappa = \sqrt {\frac {4 T}{\mathcal {N} J S ^ {2}}}.
\tag{18.13}
\end{equation*}

为研究自旋关联，限于 $N = 2$，并引入 $z$ 方向磁场，它以下列项耦合：

\begin{align*}
H ^ {\mathrm{MF}} & \rightarrow H ^ {\mathrm{MF}} - h \sum_ {i} S _ {i} ^ {z}\\
& = H ^ {\mathrm{MF}} - h \sum_{i, \, s = \pm \frac {1}{2}} s \, a _ {is} ^ {\dagger} a _ {is} .
\tag{18.14}
\end{align*}

场劈裂两个 Schwinger 玻色子色散的简并：

\begin{equation*}
\epsilon_ {\mathbf {k} m} = \epsilon_ {\mathbf {k}} - h s.
\tag{18.15}
\end{equation*}

约束方程给出

\begin{align*}
S & = \frac {1}{2} \lim_{T \to 0} \mathcal {N} ^ {- 1} \sum_ {\mathbf {k}} \left[ n \left( \epsilon_{\mathbf {k}, + \frac {1}{2}} \right) + n \left( \epsilon_{\mathbf {k}, - \frac {1}{2}} \right) \right]\\
& \rightarrow \frac {1}{2} \mathcal {N} ^ {- 1} n \left( \epsilon_{0, + \frac {1}{2}} \right) ,
\tag{18.16}
\end{align*}

磁化为\footnote{见 (6.3)。}

\begin{align*}
m _ {0} & = \frac {1}{2} \mathcal {N} ^ {- 1} \lim_{h \to 0 ^ {+}} \sum_ {\mathbf {k}} \left[ n \left( \epsilon_{\mathbf {k}, + \frac {1}{2}} \right) - n \left( \epsilon_{\mathbf {k}, - \frac {1}{2}} \right) \right]\\
& = \frac {1}{2} \mathcal {N} ^ {- 1} n \left( \epsilon_{0, + \frac {1}{2}} \right) = S.
\tag{18.17}
\end{align*}

于是，在无穷小磁场下，$T=0$ 的 Bose 凝聚发生在 $k=0$、$s=+\frac{1}{2}$。Bose 凝聚是完全的：全部可用的 Schwinger 玻色子密度（每格点 $2S$）都积聚在这一单个模式中。Schwinger 玻色子凝聚体的序参量为

\begin{equation*}
\langle a _ {\mathbf {k} s} \rangle = \langle a _ {\mathbf {k} s} ^ {\dagger} \rangle = \sqrt {2 \mathcal {N} m _ {0}} \; \delta_{s, \frac {1}{2}} \delta_{\mathbf {k}, 0}.
\tag{18.18}
\end{equation*}

这一序参量违反约束（即连接不同自旋的态）。但在平均场理论与 $1/N$ 展开之内它有精确的含义\footnote{见 Hirsch 与 Tang。}。

全部玻色子凝聚到最低态，翻译成自旋态即

\begin{equation*}
\Psi^ {\mathrm{MF}} = \prod_ {i} \left( a _ {i, + \frac {1}{2}} ^ {\dagger} \right) ^ {2S} \left| 0 \right\rangle = \left| S, S, \ldots \right\rangle ,
\tag{18.19}
\end{equation*}

与 Marshall 定理推论 5.6 给出的精确基态一致；见 (5.37)。

注意 SBMFT 色散 $\epsilon_{k}$ 在 $T \rightarrow 0$ 时约化为铁磁自旋波，按（见 (11.49)）

\begin{align*}
\lim_{T \to 0} \epsilon_ {\mathbf {k}} \rightarrow \omega_ {\mathbf {k}} & = z J S \left( 1 - \gamma_ {\mathbf {k}} \right)\\
& \sim J S | \mathbf {k} | ^ {2} + \mathcal {O} \left( | \mathbf {k} | ^ {4} \right)
\tag{18.20}
\end{align*}

消失。$\mathbf{k} = 0$ 处 Schwinger 玻色子能量的消失由 Goldstone 定理预期；见 9.2 节。

磁化率为

\begin{align*}
\chi^ {\mathrm{MF}} & = - \mathcal {N} ^ {- 1} \frac {d ^ {2} F ( h )}{d ^ {2} h} = - \frac {1}{4 \mathcal {N}} \frac {d ^ {2} F ( h )}{d ^ {2} \lambda}\\
& = \frac {1}{4 \mathcal {N} T} \sum_ {\mathbf {k}} n _ {\mathbf {k}} \left( n _ {\mathbf {k}} + 1 \right) .
\tag{18.21}
\end{align*}

平均场自旋关联由自由玻色子表达式给出：

\begin{align*}
\langle S _ {i} ^ {+} S _ {j} ^ {-} \rangle_{\mathrm{MF}} & = \langle a _ {i, + \frac {1}{2}} ^ {\dagger} a_{i, - \frac {1}{2}} a_{j, - \frac {1}{2}} ^ {\dagger} a_{j, + \frac {1}{2}} \rangle\\
& = \left| R _ {ij} \right| ^ {2} + S \delta_{ij} ,\\
R _ {ij} & = \frac {1}{\mathcal {N}} \sum_ {\mathbf {k}} n _ {\mathbf {k}} \, e ^ {i \mathbf {k} \cdot \mathbf {x} _ {ij}} ,
\tag{18.22}
\end{align*}

这里用了恒等式 (18.6)。由 (18.6) 与 (18.22)，局域自旋涨落为

\begin{equation*}
\langle S _ {i} ^ {+} S _ {i} ^ {-} \rangle_{\mathrm{MF}} = \frac {1}{\mathcal {N}} \sum_{\mathbf {k}, \mathbf {k} '} n _ {\mathbf {k}} \left( n _ {\mathbf {k} '} + 1 \right) = S ( S + 1 ).
\tag{18.23}
\end{equation*}

对 $SU(2)$，精确结果当然是 $\frac{2}{3}S(S+1)$。由求和规则 (17.43)，我们验证 (18.23) 恰好给出涨落中正确的 $(1/N)^{0}$ 阶贡献。

\subsection{一维}

一维低温 $T \ll JS$ 时，约束方程 (18.6) 中的求和由小能量、小动量区域主导。小 $k$ 时

\begin{align*}
\epsilon & \approx J S \left( \frac {1}{4} \kappa^ {2} + k ^ {2} + \dots \right) ,\\
n ( \epsilon ) & \approx T / \epsilon .
\tag{18.24}
\end{align*}

低温下把动量上截断换成无穷，得到约束求和的领头阶温度依赖：

\begin{align*}
\mathcal {N} ^ {- 1} \sum_ {\mathbf {k}} n _ {\mathbf {k}} & \sim \frac {T}{J S \pi} \int_ {0} ^ {\infty} d k \, \frac {1}{\frac {1}{4} \kappa^ {2} + k ^ {2}} + \mathcal {O} ( T ^ {2} )\\
& \sim T / ( J S \kappa ) \approx S.
\tag{18.25}
\end{align*}

于是

\begin{equation*}
\kappa \sim T / ( J S ^ {2} ) + \mathcal {O} ( T ^ {2} ).
\tag{18.26}
\end{equation*}

低温磁化率同样由积分的小动量部分主导。把 (18.21) 的项按 $k$ 幂次展开，可算出其领头温度依赖：

\begin{align*}
\chi^ {\mathrm{MF}} & = \frac {1}{2 T} \int_{- \pi} ^ {\pi} \frac {d k}{2 \pi} \, n _ {k} \left( n _ {k} + 1 \right)\\
& \sim \frac {J S ^ {4}}{T ^ {2}}.
\tag{18.27}
\end{align*}

$S = \frac{1}{2}$ Heisenberg 铁磁体的精确磁化率 $\chi$ 已由 Takahashi 用 Bethe 解算出：

\begin{equation*}
\chi^ {S = \frac {1}{2}} \sim \frac {J}{24} T ^ {- 2} \sim \frac {2}{3} \chi^ {\mathrm{MF}}
\tag{18.28}
\end{equation*}

这一 $2/3$ 因子的偏差先前已在局域自旋涨落 (18.23) 中出现。然而在 $\chi$ 这样的长程自旋关联函数中也发现同一因子，令人惊讶。

(18.22) 的 $R_{ij}$ 在大距离下可由被积函数在低动量的行为计算：

\begin{equation*}
R _ {ij} \sim \frac {T}{J S \pi} \int_ {0} ^ {\infty} d k \, \frac {e ^ {i k | i - j |}}{\frac {1}{4} \kappa^ {2} + k ^ {2}} = S \, e ^ {- \frac {1}{2} \kappa | i - j |} ,
\tag{18.29}
\end{equation*}

这里用了 (18.26)。低温大距离下用 (18.6)：

\begin{equation*}
\langle S _ {i} ^ {z} S _ {j} ^ {z} \rangle_{\mathrm{MF}} \sim \frac {1}{2} S ^ {2} e ^ {- | i - j | / \xi} ,
\tag{18.30}
\end{equation*}

确立参数 $\kappa$ 是逆关联长度：

\begin{equation*}
\xi = \kappa^ {- 1} \sim J S ^ {2} / T.
\tag{18.31}
\end{equation*}

这一结果与 M. Fisher 计算的一维经典 Heisenberg 铁磁体精确关联长度一致。

\subsection{二维}

考虑正方格子，$z = 4$。计算动量求和的有用函数是

\begin{equation*}
\rho ( \gamma ) = \mathcal {N} ^ {- 1} \sum_ {\mathbf {k}} \delta \left( \gamma - \gamma_ {\mathbf {k}} \right) = \frac {2}{\pi^ {2}} K \left( 1 - \gamma^ {2} \right) ,
\tag{18.32}
\end{equation*}

它描述正方格子上紧束缚跳跃矩阵的态密度，其中

\begin{equation*}
\gamma_ {\mathbf {k}} = \frac {1}{2} \left( \cos k _ {x} + \cos k _ {y} \right).
\tag{18.33}
\end{equation*}

$K(m)$ 是“第一类完全椭圆积分”\footnote{见 Abramowitz 与 Stegun，p. 590。}：

\begin{align*}
K ( m ) & = \int_ {0} ^ {1} d t \left[ \left( 1 - t ^ {2} \right) \left( 1 - m t ^ {2} \right) \right] ^ {- \frac {1}{2}}\\
& \approx \left\{ \begin{array}{ll} \theta ( m ) \left[ \frac {\pi}{2} + \frac {\pi}{8} m + \mathcal {O} ( m ) ^ {2} \right] & m \ll 1 \\ \frac {1}{2} \ln [ 16 / ( 1 - m ) ] & | 1 - m | \ll 1 \end{array} \right. .
\tag{18.34}
\end{align*}

$\rho(\gamma)$ 的重要特征是 $\gamma = \pm 1$ 处的阶跃不连续与 $\gamma_{k} = 0$ 附近的对数发散。该函数画于图~\ref{fig:18.2}。由于主导的低温关联依赖 $\gamma \approx 1$ 处的 $\rho (\gamma)$，用替换

\begin{figure}[H]
\centering
\includegraphics[width=0.55\textwidth]{fig18_2.jpg}
\caption{正方格子上紧束缚色散 $\gamma_{k}$ 的态密度。}
\label{fig:18.2}
\end{figure}

\begin{align*}
\rho ( \gamma ) & \rightarrow \bar {\rho} = \frac {1}{\pi} \, , \quad \gamma \in [ 1 - \pi, 1 ] ,\\
\int_{1 - \pi} ^ {1} d \gamma \, \bar {\rho} ( \gamma ) & = 1 ,
\tag{18.35}
\end{align*}

简化后续计算——延长 $\gamma$ 的定义域以保持单位总密度。用这一相当无害的近似，约束方程容易求解：

\begin{align*}
\mathcal {N} ^ {- 1} \sum_ {\mathbf {k}} n _ {\mathbf {k}} & \approx \frac {T}{4 J S \pi} \int_ {1 - \pi} ^ {1} d \gamma \, \frac {1}{4 ( 1 - \gamma ) + \kappa^ {2} / 4}\\
& = \frac {T}{J S \pi} \ln \left( 16 \pi / \kappa^ {2} \right) + \mathcal {O} ( T ^ {2} ) = S.
\tag{18.36}
\end{align*}

反解此式并用 (18.11)，得领头阶温度解：

\begin{equation*}
\kappa \approx \sqrt {16 \pi} \, \exp \left( \frac {- 2 \pi J S ^ {2}}{T} \right).
\tag{18.37}
\end{equation*}

大距离的自旋关联可用 $R_{ij}$ 的 Ornstein--Zernicke 近似\footnote{见 (13.48)。}计算：

\begin{align*}
R _ {ij} & \approx \frac {T}{J S} \int \frac {d ^ {2} \mathbf {k}}{( 2 \pi ) ^ {2}} \frac {e ^ {i \mathbf {k} \mathbf {x} _ {ij}}}{\frac {1}{4} \kappa^ {2} + k ^ {2}}\\
& \propto \left( | \mathbf {x} _ {ij} | / \xi \right) ^ {- \frac {1}{2}} \exp \left( - | \mathbf {x} _ {ij} | \kappa / 2 \right) \left( 1 + \frac {2}{| \mathbf {x} _ {ij} | \kappa} \dots \right) ,
\tag{18.38}
\end{align*}

$\mathbf{x}_{ij} = \mathbf{x}_i - \mathbf{x}_j$。由 (18.22)：

\begin{equation*}
\langle S _ {i} ^ {z} S _ {j} ^ {z} \rangle_{\mathrm{MF}} \propto \frac {\xi}{| \mathbf {x} _ {ij} |} \, e ^ {- | \mathbf {x} _ {ij} | / \xi} ,
\tag{18.39}
\end{equation*}

其中

\begin{equation*}
\xi = \kappa^ {- 1} \, ,
\tag{18.40}
\end{equation*}

即 $\kappa^{-1}$ 是自旋关联长度。于是我们得出结论：关联长度随温度降低指数增长，且任何有限 $T > 0$ 没有长程序——与 Mermin--Wagner 定理一致。

回忆 13.1 节曾看到，经典 Heisenberg 铁磁体在长距离低温下由耦合常数为

\begin{equation*}
f = \frac {T}{J S ^ {2}}.
\tag{18.41}
\end{equation*}

的非线性 {$\sigma$} 模型描述。(18.37) 恢复了到单圈阶的重整化群结果 (13.55) 与 $CP^{N-1}$ 模型的大 $N$ 近似 (14.26)。

\section{反铁磁体情形}

由 (16.31) 与 (18.2)，AFM-B 表示的平均场哈密顿量为

\begin{align*}
H ^ {\mathrm{MF}} & = \sum_{i, m} \lambda \, a _ {im} ^ {\dagger} a _ {im} + Q \sum_{\langle ij \rangle, m} \left( a _ {im} ^ {\dagger} a _ {jm} ^ {\dagger} + a _ {im} a _ {jm} \right)\\
& \quad + N \mathcal {N} \frac {z Q ^ {2}}{2 J} - N \mathcal {N} S \lambda\\
& = \sum_{\mathbf {k} m} \left[ \lambda \, a _ {\mathbf {k} m} ^ {\dagger} a _ {\mathbf {k} m} + \frac {1}{2} z Q \gamma_ {\mathbf {k}} \left( a _ {\mathbf {k} m} ^ {\dagger} a _ {- \mathbf {k} m} ^ {\dagger} + a _ {\mathbf {k} m} a _ {- \mathbf {k} m} \right) \right]\\
& \quad + N \mathcal {N} \frac {z Q ^ {2}}{2 J} - N \mathcal {N} S \lambda .
\tag{18.42}
\end{align*}

对 $N = 2$，SBMFT 哈密顿量类似于 (11.60) 的 Holstein--Primakoff 自旋波哈密顿量 $H_{1}$，只是这里两个 Schwinger 玻色子味替换了单个 Holstein--Primakoff 玻色子。与自旋波问题极类似，$H^{\mathrm{MF}}$ 可用正则 Bogoliubov 变换对角化：

\begin{equation*}
\alpha_ {\mathbf {k} m} = \cosh \theta_ {\mathbf {k}} \, a _ {\mathbf {k} m} - \sinh \theta_ {\mathbf {k}} \, a _ {- \mathbf {k} m} ^ {\dagger} ,
\tag{18.43}
\end{equation*}

或反解

\begin{equation*}
a _ {\mathbf {k} m} = \cosh \theta_ {\mathbf {k}} \, \alpha_ {\mathbf {k} m} + \sinh \theta_ {\mathbf {k}} \, \alpha_{- \mathbf {k} m} ^ {\dagger}.
\tag{18.44}
\end{equation*}

把 (18.44) 代入 (18.42)，得到 $\alpha$ 玻色子的正规对角哈密顿量：

\begin{align*}
H ^ {\mathrm{MF}} & = \frac {1}{2} \sum_{\mathbf {k} m} \Big[ \left( \lambda \cosh 2 \theta_ {\mathbf {k}} + z Q \gamma_ {\mathbf {k}} \sinh 2 \theta_ {\mathbf {k}} \right) \left( \alpha_{\mathbf {k} m} ^ {\dagger} \alpha_{\mathbf {k} m} + \alpha_{\mathbf {k} m} \alpha_{\mathbf {k} m} ^ {\dagger} \right)\\
& \quad + \left( \lambda \sinh 2 \theta_ {\mathbf {k}} + z Q \gamma_ {\mathbf {k}} \cosh 2 \theta_ {\mathbf {k}} \right) \left( \alpha_{\mathbf {k} m} ^ {\dagger} \alpha_{- \mathbf {k} m} ^ {\dagger} + \alpha_{\mathbf {k} m} \alpha_{- \mathbf {k} m} \right) \Big]\\
& \quad + N \mathcal {N} \frac {z Q ^ {2}}{2 J} - \mathcal {N} N \left( S + \frac {1}{2} \right) \lambda .
\tag{18.45}
\end{align*}

选 $\theta_{\mathbf{k}}$ 使反常项 $\alpha^{\dagger}\alpha^{\dagger}$ 与 $\alpha \alpha$ 消失，即对 $\theta_{\mathbf{k}}$ 的条件：

\begin{equation*}
\tanh 2 \theta_ {\mathbf {k}} = - \frac {z Q \gamma_ {\mathbf {k}}}{\lambda}.
\tag{18.46}
\end{equation*}

解出 $\theta_{k}$ 后，把 (18.45) 中的双曲函数替换为 (18.46) 右端的有理函数，得到正规对角哈密顿量：

\begin{align*}
H ^ {\mathrm{MF}} & = \sum_{\mathbf {k} m} \omega_ {\mathbf {k}} \left( \alpha_{\mathbf {k} m} ^ {\dagger} \alpha_{\mathbf {k} m} + \frac {1}{2} \right) + N \mathcal {N} \frac {z Q ^ {2}}{2 J} - N \mathcal {N} \left( S + \frac {1}{2} \right) \lambda ,\\
\omega_ {\mathbf {k}} & = \sqrt {\lambda^ {2} - \left( z Q \gamma_ {\mathbf {k}} \right) ^ {2}} .
\tag{18.47}
\end{align*}

平均场自由能为

\begin{equation*}
F ^ {\mathrm{MF}} = \beta^ {- 1} \sum_{\mathbf {k} m} \ln \left[ 2 \sinh \left( \frac {\beta \omega_ {\mathbf {k}}}{2} \right) \right] - \mathcal {N} N \left( S + \frac {1}{2} \right) \lambda + N \mathcal {N} \frac {z Q ^ {2}}{2 J}.
\tag{18.48}
\end{equation*}

平均场方程由 (18.48) 对 $\lambda$ 与 $Q$ 求导：

\begin{equation*}
\frac {1}{\mathcal {N}} \sum_ {\mathbf {k}} \frac {\lambda}{\sqrt {\lambda^ {2} - \left( z Q \gamma_ {\mathbf {k}} \right) ^ {2}}} \left( n _ {\mathbf {k}} + \frac {1}{2} \right) = S + \frac {1}{2} ,
\tag{18.49}
\end{equation*}

\begin{equation*}
\frac {1}{\mathcal {N}} \sum_{\mathbf {k}} \frac {z ^ {2} \gamma_{\mathbf {k}} ^ {2} Q}{\sqrt {\lambda^ {2} - \left( z Q \gamma_{\mathbf {k}} \right) ^ {2}}} \left( n _ {\mathbf {k}} + \frac {1}{2} \right) = \frac {z Q}{J}.
\tag{18.50}
\end{equation*}

平均场基态 $\Psi^{\mathrm{MF}}$ 是所有 $\alpha$ 的真空：

\begin{equation*}
\alpha_{\mathbf {k}, m} \Psi_ {0} ^ {\mathrm{MF}} = 0 \, , \quad \forall \mathbf {k}, m.
\tag{18.51}
\end{equation*}

用 (18.43)，可把 $\Psi^{\mathrm{MF}}$ 用原 Schwinger 玻色子显式写出：

\begin{align*}
\Psi^ {\mathrm{MF}} & = C \exp \left[ \frac {1}{2} \sum_ {ij} u _ {ij} \left( \sum_ {m} a _ {im} ^ {\dagger} a _ {jm} ^ {\dagger} \right) \right] \left| 0 \right\rangle ,\\
u _ {ij} & = \frac {1}{\mathcal {N}} \sum_ {\mathbf {k}} e ^ {i \mathbf {k} \mathbf {x} _ {ij}} \tanh \theta_ {\mathbf {k}}.
\tag{18.52}
\end{align*}

$N = 2$ 时，用未转动算符 $a^{\dagger}, b^{\dagger}$，$\Psi^{\mathrm{MF}}$ 就是 Schwinger 玻色子平均场态 (8.5)：

\begin{equation*}
\Psi_{N = 2} ^ {\mathrm{MF}} = \left| \hat {u} \right\rangle = \exp \left[ \sum_{i \in A, \, j \in B} u _ {ij} \left( a _ {i} ^ {\dagger} b _ {j} ^ {\dagger} - b _ {i} ^ {\dagger} a _ {j} ^ {\dagger} \right) \right] \left| 0 \right\rangle .
\tag{18.53}
\end{equation*}

$\Psi^{\mathrm{MF}}$ 含许多占据数不同于 $2S$ 的组态，因此不是纯自旋态。如第 8 章所示，经 Gutzwiller 投影后它化为价键态。由于

\begin{equation*}
\tanh \left( \theta_{\mathbf {k} + \vec {\pi}} \right) = - \tanh \left( \theta_{\mathbf {k}} \right) ,
\tag{18.54}
\end{equation*}

$\vec{\pi} = (\pi, \pi, \ldots)$，键参数 $u_{ij}$ 只连接子晶格 $A$ 与 $B$。而且可以验证，对上述近邻模型 $u_{ij} \geq 0$，故价键态满足 (5.13) 定义的 Marshall 符号。

虽然 $\Psi^{\mathrm{MF}}$ 明显转动不变，却可能有或没有长程反铁磁序——这取决于 $u_{ij}$ 的长距离衰减。我们将看到，近邻模型的 SBMFT 基态在一维无序、二维有长程序。

为进一步计算，引入参数化：

\begin{align*}
\omega_ {\mathbf {k}} & \equiv c \sqrt { ( \kappa / 2 ) ^ {2} + \frac {z}{2} \left( 1 - \gamma_{\mathbf {k}} ^ {2} \right) } ,\\
c & \equiv Q \sqrt {2 z} ,\\
\kappa & \equiv \frac {2}{c} \sqrt {\lambda^ {2} - \left( z Q \right) ^ {2}} ,\\
t & = \frac {T}{z Q}.
\tag{18.55}
\end{align*}

$c$、$\kappa$、$t$ 分别描述自旋波速度、逆关联长度与无量纲温度。图~\ref{fig:18.3} 画出一维反铁磁体的色散。由 (18.48)，在区中心与区角附近，平均场色散是自由有质量相对论性玻色子：

\begin{equation*}
\omega_{\mathbf {k}} ^ {\gamma} \approx c \sqrt { ( \kappa / 2 ) ^ {2} + \left| \mathbf {k} - \mathbf {k} _ {\gamma} \right| ^ {2} } \, , \quad \mathbf {k} _ {\gamma} = 0, \vec {\pi}.
\tag{18.56}
\end{equation*}

当能隙（或“质量” $c\kappa/2$）消失时，$\omega_{k}^{\alpha}$ 是 Goldstone 模，约化为反铁磁自旋波色散 (11.65)。

\begin{figure}[H]
\centering
\includegraphics[width=0.55\textwidth]{fig18_3.jpg}
\caption{一维反铁磁体的平均场色散。}
\label{fig:18.3}
\end{figure}

自旋关联函数由在 (16.40) 中用 (18.44) 把 $a$ 算符换成 $\alpha$ 算符得到：

\begin{align*}
S ^ {\mathrm{MF}} ( \mathbf {q} ) & = \frac {1}{\mathcal {N}} \langle S _ {\mathbf {q}} ^ {+} S _ {- \mathbf {q}} ^ {-} \rangle_{\mathrm{MF}}\\
& = \frac {1}{\mathcal {N}} \sum_ {\mathbf {k}} \Bigg\{ \cosh \left[ 2 \left( \theta_{\mathbf {k}} + \theta_{\mathbf {k} + \mathbf {q} + \vec {\pi}} \right) \right]\\
& \quad \times \left( n_{\mathbf {k}} + \frac {1}{2} \right) \left( n_{\mathbf {k} + \mathbf {q} + \vec {\pi}} + \frac {1}{2} \right) - \frac {1}{4} \Bigg\} .
\tag{18.57}
\end{align*}

用 (18.49) 确认求和规则 (17.43) 的大 $N$ 极限：

\begin{equation*}
\frac {1}{\mathcal {N}} \sum_ {\mathbf {q}} S ^ {\mathrm{MF}} ( \mathbf {q} ) = S ( S + 1 ).
\tag{18.58}
\end{equation*}

$N = 2$ 时，平均场求和规则超出精确结果一个熟悉的因子 $\frac{3}{2}$。

自旋关联在 $x_{ij} = x_i - x_j$ 处的空间依赖为

\begin{equation*}
S ^ {\mathrm{MF}} ( \mathbf {x} _ {ij} ) = \left| f ( \mathbf {x} _ {ij} ) \right| ^ {2} - \left| g ( \mathbf {x} _ {ij} ) \right| ^ {2} - \frac {1}{4} \delta_{ij} ,
\tag{18.59}
\end{equation*}

低温长距离下

\begin{align*}
f ( \mathbf {x} _ {ij} ) & = \mathcal {N} ^ {- 1} \sum_ {\mathbf {k}} \frac {\left( n _ {\mathbf {k}} + \frac {1}{2} \right) e ^ {i \mathbf {k} \mathbf {x} _ {ij}}}{\sqrt {1 - \gamma_{\mathbf {k}} ^ {2} \left[ \kappa^ {2} / (2z) + 1 \right] ^ {- 1}}}\\
& \approx 2 z t \left( 1 + e ^ {i \vec {\pi} \mathbf {X} _ {ij}} \right) \int \frac {d ^ {d} \mathbf {k}}{( 2 \pi ) ^ {d}} \frac {e ^ {i \mathbf {k} \mathbf {x} _ {ij}}}{( \kappa / 2 ) ^ {2} + | \mathbf {k} | ^ {2}}\\
& \propto \left( 1 + e ^ {i \vec {\pi} \mathbf {X} _ {ij}} \right) \left( | \mathbf {x} _ {ij} | / \xi \right) ^ {- ( d - 1 )/2} \exp \left( - | \mathbf {x} _ {ij} | \kappa / 2 \right) ,
\tag{18.60}
\end{align*}

小 $\kappa$ 低温下的长距离行为用了 (13.48)。类似地

\begin{align*}
g ( \mathbf {x} _ {ij} ) & = \mathcal {N} ^ {- 1} \sum_ {\mathbf {k}} \gamma_ {\mathbf {k}} \frac {\left( n _ {\mathbf {k}} + \frac {1}{2} \right) e ^ {i \mathbf {k} \mathbf {x} _ {ij}}}{\sqrt {1 - \gamma_{\mathbf {k}} ^ {2} \left[ \kappa^ {2} / (2z) + 1 \right] ^ {- 1}}}\\
& \propto \left( 1 - e ^ {i \vec {\pi} \mathbf {x} _ {ij}} \right) \left( | \mathbf {x} _ {ij} | / \xi \right) ^ {- ( d - 1 )/2} \exp \left( - | \mathbf {x} _ {ij} | \kappa / 2 \right)
\tag{18.61}
\end{align*}

于是对 $\kappa > 0$：

\begin{equation*}
S ^ {\mathrm{MF}} ( \mathbf {x} _ {ij} ) \propto e ^ {i \vec {\pi} \mathbf {x} _ {ij}} \left( \frac {\xi}{| \mathbf {x} _ {ij} |} \right) ^ {d - 1} \exp \left( - | \mathbf {x} _ {ij} | / \xi \right) ,
\tag{18.62}
\end{equation*}

关联长度为 $\xi = \kappa^{-1}$。

均匀磁化率直接由 (18.57) 用恒等式

\begin{equation*}
\chi_ {0} ^ {\mathrm{MF}} = \frac {1}{2 T} S ^ {\mathrm{MF}} ( \mathbf {q} = 0 ) = \frac {1}{T} \sum_ {\mathbf {k}} n _ {\mathbf {k}} \left( n _ {\mathbf {k}} + 1 \right).
\tag{18.63}
\end{equation*}

得到。

\subsection{长程反铁磁序}

没有任何磁场时，基态是单态。在二维及更高维，第 11 章我们看到近邻反铁磁体对所有 $S \geq \frac{1}{2}$ 有长程序。为考察自发破缺对称的可能性，引入耦合到交错磁化的无穷小序场 $h$。限于 $N = 2$：

\begin{align*}
H ^ {\mathrm{MF}} & \rightarrow H ^ {\mathrm{MF}} - h \sum_ {i} S _ {i} ^ {z}\\
& = H ^ {\mathrm{MF}} - h \sum_{i, \, s = - \frac {1}{2}, \frac {1}{2}} s \, a _ {is} ^ {\dagger} a _ {is} ,
\tag{18.64}
\end{align*}

回忆 Schwinger 玻色子在 (16.11) 中是用子晶格转动后的自旋方向定义的。(18.64) 可用依赖自旋的变换角 $\theta_{k,s}$ 对角化。重复导出 (18.47) 的步骤：

\begin{equation*}
\omega_{\mathbf {k}, s} = \sqrt { \left( \lambda - s h \right) ^ {2} - \left( z Q \gamma_{\mathbf {k}} \right) ^ {2} }.
\tag{18.65}
\end{equation*}

自发交错磁化由极限

\begin{align*}
m _ {0} & = \lim_{h \to 0 ^ {+}} m ( h ) \, ,\\
m ( h ) & = \lim_{\mathcal {N} \to \infty} \frac {1}{\mathcal {N}} \sum_{i, s} s \left\langle a _ {is} ^ {\dagger} a _ {is} \right\rangle_{h}\\
& = \lim_{\mathcal {N} \to \infty} \frac {1}{\mathcal {N}} \sum_{\mathbf {k} s} \frac {\lambda - s h}{\sqrt { \left( \lambda - s h \right) ^ {2} - \left( z Q \gamma_{\mathbf {k}} \right) ^ {2} }} \left[ n \left( \omega_{\mathbf {k}, s} \right) + \frac {1}{2} \right]
\tag{18.66}
\end{align*}

给出。约束方程 (18.49) 为

\begin{equation*}
\frac {1}{2 \mathcal {N}} \sum_{\mathbf {k}, s} \frac {\lambda - s h}{\sqrt { \left( \lambda - s h \right) ^ {2} - \left( z Q \gamma_{\mathbf {k}} \right) ^ {2} }} \left[ n \left( \omega_{\mathbf {k}, s} \right) + \frac {1}{2} \right] = S + \frac {1}{2}.
\tag{18.67}
\end{equation*}

再用 $c, \kappa, \tilde{h}$ 参数化色散：

\begin{align*}
\omega_{\mathbf {k}, s} & \equiv c \sqrt { \kappa^ {2} /4 - \left( s - \frac {1}{2} \right) z \tilde {h} + \frac {z}{2} \left( 1 - \gamma_{\mathbf {k}} ^ {2} \right) + \mathcal {O} ( \tilde {h} ^ {2} ) } ,\\
c & = \sqrt {2 z} \, Q ,\\
\kappa & = \frac {2}{c} \sqrt { \left( \lambda - \frac {1}{2} h \right) ^ {2} - \left( z Q \right) ^ {2} } ,\\
\tilde {h} & = h \frac {\lambda}{\left( z Q \right) ^ {2}} .
\tag{18.68}
\end{align*}

若

\begin{equation*}
\lim_{\mathcal {N} \to \infty} \kappa > 0 ,
\tag{18.69}
\end{equation*}

则 (18.66) 中 $s = \pm \frac{1}{2}$ 的两项都在 $h = 0$ 连续，因此

\begin{equation*}
\lim_{h \to 0 ^ {+}} m ( h, \kappa ) = 0 ,
\tag{18.70}
\end{equation*}

即无自发对称破缺。反之，若

\begin{equation*}
\kappa = \mathcal {O} ( \mathcal {N} ) ^ {- 1}
\tag{18.71}
\end{equation*}

则 $k = 0, \vec{\pi}$ 处 $s = +\frac{1}{2}$ 的项贡献 $N$ 阶的量，对交错磁化 (18.66) 给出宏观贡献（即一阶的量）。由于它同时也代表对 Schwinger 玻色子密度 (18.67) 的宏观贡献，可以说 $s = +\frac{1}{2}$ 玻色子在 $k = 0, \vec{\pi}$ 处发生 Bose 凝聚（见 (18.18) 之后的讨论）。凝聚体的序参量于是为

\begin{equation*}
\langle a _ {\mathbf {k} s} \rangle = \langle a _ {\mathbf {k} s} ^ {\dagger} \rangle = \sqrt {\frac {\mathcal {N} m _ {0}}{2}} \; \delta_{s, \frac {1}{2}} \left( \delta_{\mathbf {k}, 0} + \delta_{\mathbf {k}, \vec {\pi}} \right).
\tag{18.72}
\end{equation*}

为计算 $m_0 = \lim_{h\to 0}m(h)$，从 (18.67) 减去 (18.66)，消去发散的 $s = \frac{1}{2}$ 项：

\begin{align*}
m _ {0} & = S + \frac {1}{2}\\
& \quad - \lim_{h \to 0 ^ {+}} \lim_{\mathcal {N} \to \infty} \mathcal {N} ^ {- 1} \sum_ {\mathbf {k}} \frac {2 + 2 \tilde {h} + \kappa^ {2} /4}{\sqrt {2 \tilde {h} + \kappa^ {2} /4 + 2 \left( 1 - \gamma_{\mathbf {k}} ^ {2} \right) }} \left[ n \left( \omega_{\mathbf {k} \frac {1}{2}} \right) + \frac {1}{2} \right].
\tag{18.73}
\end{align*}

保持 $h > 0$ 使谱与被积项分母中有能隙。于是在热力学极限下允许把 (18.73) 换成积分；$T = 0$ 时取 $n(\omega_{\mathbf{k}, + \frac{1}{2}}) = 0$ 与 $\kappa = 0$：

\begin{equation*}
m _ {0} = S + \frac {1}{2} - \frac {1}{2} \int \frac {d ^ {d} \mathbf {k}}{( 2 \pi ) ^ {d}} \frac {1}{\sqrt {1 - \gamma_{\mathbf {k}} ^ {2}}}.
\tag{18.74}
\end{equation*}

对 $d$ 维立方格子，积分给出数值结果

\begin{equation*}
m _ {0} ( d ) = \left\{ \begin{array}{ll} 0 & d = 1 \\ S - 0.19660 & d = 2 \\ S - 0.078 & d = 3 \end{array} \right..
\tag{18.75}
\end{equation*}

注意，与铁磁情形 (18.17) 相反，有序矩总是小于经典值 $S$。这源于量子零点运动——其根源是哈密顿量与交错磁化不可对易。SBMFT 的 $m_{0}$ 结果与 (11.69) 的低阶自旋波理论一致。

\subsection{一维}

零温时取 $n_{\mathbf{k}} = 0$，展开 (18.49)：

\begin{align*}
S + \frac {1}{2} & = \frac {1}{2} \int_{- \pi} ^ {\pi} \frac {d k}{2 \pi} \frac {1}{\sqrt {1 - \frac {1}{1 + \kappa^ {2} /8} \cos^ {2} ( k ) }}\\
& = \frac {1}{\pi} K \left( \frac {1}{1 + \kappa^ {2} /8} \right)\\
& \sim \frac {1}{\pi} \ln \left( \frac {8 \sqrt {2}}{\kappa} \right) + \mathcal {O} ( \kappa ) ,
\tag{18.76}
\end{align*}

得到

\begin{equation*}
\kappa \approx \sqrt {32} \, \exp \left[ - \pi \left( S + \frac {1}{2} \right) \right].
\tag{18.77}
\end{equation*}

(18.75) 中已发现一维反铁磁体的 SBMFT 基态不可能有长程序。由于 $\kappa$ 随 $S$ 指数下降，可以像忽略 $S^{-1}$ 更高阶修正那样忽略 $\kappa$。从 (18.49) 减去 (18.50)：

\begin{align*}
c & = J \left( S + \frac {1}{2} - \frac {1}{2} \int_{- \pi} ^ {\pi} \frac {d k}{2 \pi} \frac {\sin^ {2} k}{\sqrt {\kappa^ {2} /8 + \sin^ {2} k}} \right)\\
& \approx J \left[ S + \frac {1}{2} - \frac {2}{\pi} + \mathcal {O} ( \kappa, S ^ {- 1} ) \right] .
\tag{18.78}
\end{align*}

虽然 $c$ 与其经典值 $c = JS$（见 (11.41)）相去不远，基态关联却按指数衰减。(18.77) 给出的关联长度 $\kappa^{-1}$ 与连续途径 (15.14) 给出的关联长度一致。

平均场激发并非 Heisenberg 模型的物理激发（例如它们包含违反约束的电荷涨落）。但其谱中的能隙 $c\kappa/2$ 与一维整数自旋链 Haldane 能隙的存在一致。物理磁振子谱可从 $\operatorname{Im} S(\mathbf{q},\omega)$ 的峰推得。由于自旋 1 的激发至少涉及两个 Schwinger 玻色子，物理磁振子具有能隙

\begin{equation*}
\Delta = c \kappa .
\tag{18.79}
\end{equation*}

于是，SBMFT 在无 $\Theta$ 项时恢复了 Haldane 的连续理论结果（见第 15 章）。

必须警惕：SBMFT 对半奇整数自旋链失效。摧毁 Haldane 能隙的 $\Theta$ 项 (15.3) 的效应显然在平均场理论中缺席。SBMFT 的非简并基态违反 5.2 节 Lieb--Schultz--Mattis 定理对半奇整数自旋的结论。Read 与 Sachdev 通过在大 $N$ 理论中引入 $\Theta$ 项克服了这一问题；他们为 $Q$、$\lambda$ 涨落发展了一个连续规范理论（见文献注释）。

\subsection{二维}

二维中，由 Mermin--Wagner 定理预期有限温无长程序。的确我们将发现，低温下平均场方程给出有限的逆关联长度 $\kappa(T,S)$。

平均场方程可解出 $\kappa(T, S), c(T, S)$。利用 $\omega_{\mathbf{k}} = \omega(\gamma_{\mathbf{k}})$ 并替换

\begin{equation*}
\frac {1}{\mathcal {N}} \sum_ {\mathbf {k}} F ( \omega_{\mathbf {k}} ) \rightarrow 2 \int_ {0} ^ {1} d \gamma \, \rho ( \gamma ) F [ \omega ( \gamma ) ] ,
\tag{18.80}
\end{equation*}

是方便的，正方格子的 $\rho(\gamma)$ 为

\begin{equation*}
\rho ( \gamma ) = \frac {2}{\pi^ {2}} K \left( 1 - \gamma^ {2} \right).
\tag{18.81}
\end{equation*}

结果表明，温度高于 $T_{max} > 0.91J$ 时，平均场方程无非平凡（$Q \neq 0$）解\footnote{见 Arovas 与 Auerbach。}。这反映了 SBMFT 无法描述高温无序相——那里近邻关联已被摧毁。在关联长度大的温度，即

\begin{equation*}
\kappa \ll t \ll 1 ,
\tag{18.82}
\end{equation*}

约束方程可循 Takahashi 展开：

\begin{align*}
S + \frac {1}{2} & = \frac {1}{2} \int_ {- 1} ^ {1} d \gamma \, \rho ( \gamma ) \left( 1 + \kappa^ {2} /8 - \gamma^ {2} \right) ^ {- \frac {1}{2}} \coth \left[ \left( 1 + \kappa^ {2} /8 - \gamma^ {2} \right) ^ {\frac {1}{2}} / 2t \right]\\
& = \frac {t}{\pi} \left[ \log \left( \frac {32}{\kappa^ {2}} \right) - \log \left( \frac {2}{t} \right) \right]\\
& \quad + \frac {1}{\pi^ {2}} \int_ {- 1} ^ {1} d \gamma \left( 1 - \gamma^ {2} \right) ^ {- \frac {1}{2}} K \left( 1 - \gamma^ {2} \right) + \mathcal {O} ( t, \kappa ) .
\tag{18.83}
\end{align*}

类似地，从 (18.49) 减去 (18.50) 并把积分展开到 $\kappa, t$ 的低阶：

\begin{align*}
& S + \frac {1}{2} - \frac {4 Q}{J}\\
\geq \; & \frac {1}{2} \int_ {- 1} ^ {1} d \gamma \, \rho ( \gamma ) \left( 1 + \kappa^ {2} /8 - \gamma^ {2} \right) ^ {\frac {1}{2}} \coth \left[ \left( 1 + \kappa^ {2} /8 - \gamma^ {2} \right) ^ {\frac {1}{2}} / 2t \right]\\
\approx \; & \frac {1}{\pi^ {2}} \int_ {- 1} ^ {1} d \gamma \left( 1 - \gamma^ {2} \right) ^ {\frac {1}{2}} K \left( 1 - \gamma^ {2} \right) + \mathcal {O} ( t ^ {3}, \kappa ) .
\tag{18.84}
\end{align*}

由 (18.74) 与 (18.75)，有序矩为

\begin{equation*}
m _ {0} = S + \frac {1}{2} - \frac {1}{\pi^ {2}} \int_ {- 1} ^ {1} d \gamma \left( 1 - \gamma^ {2} \right) ^ {- \frac {1}{2}} K \left( 1 - \gamma^ {2} \right) = S - 0.19660.
\tag{18.85}
\end{equation*}

零温自旋波速度为

\begin{equation*}
c = \sqrt {8} \, J S \, Z _ {c} ,
\tag{18.86}
\end{equation*}

自旋波速度重整化因子由 (18.84) 得

\begin{align*}
Z _ {c} & = 1 + S ^ {- 1} \left[ \frac {1}{2} - \frac {1}{\pi^ {2}} \int_ {- 1} ^ {1} d \gamma \left( 1 - \gamma^ {2} \right) ^ {\frac {1}{2}} K \left( 1 - \gamma^ {2} \right) \right]\\
& = 1 + 0.078974 / S .
\tag{18.87}
\end{align*}

渐近自旋关联由 (18.62) 给出，关联长度为 $\xi = \kappa^{-1}$。反解 (18.83) 并用 (18.84)，得到依赖温度的关联长度

\begin{equation*}
\xi = \frac {\sqrt {2} J S Z _ {c}}{T} \exp \left[ - \frac {2 \pi Z _ {c} J S m _ {0}}{T} \right] \left[ 1 + \mathcal {O} ( t ^ {2} ) \right].
\tag{18.88}
\end{equation*}

把 (18.88) 与连续近似值\footnote{见 Chakravarty、Halperin 与 Nelson，第 15 章文献注释。}(15.17) 比较，可以如下确定正方格子模型的重整化劲度常数 $\rho_{s}(S)$：

\begin{equation*}
\rho_ {s} = \lim_{T \to 0} \frac {T}{2 \pi} \log ( \xi ) = J S m _ {0} Z _ {c} ,
\tag{18.89}
\end{equation*}

$m_{0}$ 与 $Z_{c}$ 分别由 (18.85) 与 (18.87) 给出。大 $S$ 极限下它正确恢复经典值 $\rho_{s} \rightarrow JS^{2}$。

\begin{exercises}

\item 用零温下 $Q$ 与 $\lambda$ 的显式解，证明 FM-B 与 AFM-B 平均场理论的基态能量为

\begin{equation*}
E ^ {\mathrm{MF}} = \lim_{T \to 0} F ^ {\mathrm{MF}} = - N \mathcal {N} \frac {z Q ^ {2}}{2 J}.
\tag{18.90}
\end{equation*}

加上从哈密顿量中丢掉的常数 $zNJS^{2}/2$（见 (16.7)），证明 $N = 2$ 时恢复正确的铁磁基态能量。

\item 证明反铁磁 SBMFT 的交错磁场项 (18.64) 导致 (18.65) 的色散 $\omega_{k_{s}}$。提示：求二次型矩阵

\begin{equation*}
L = \sum_ {s} \left( z_{\mathbf {k} s} ^ {\dagger}, \bar {z}_{- \mathbf {k} s} \right) \left( \begin{array}{cc} \omega - \lambda + \frac {1}{2} h _ {s} & z \gamma_{\mathbf {k}} Q \\ z \gamma_{\mathbf {k}} Q & - \omega - \lambda + \frac {1}{2} h _ {s} \end{array} \right) \binom{z_{\mathbf {k} s}}{\bar {z}_{- \mathbf {k} s} ^ {\dagger}}
\tag{18.91}
\end{equation*}

行列式的零点，其中 $z$ 与 $\bar{z}$ 分别是子晶格 A、B 的相干态变量。

\item 对反铁磁模型，给 (18.42) 的 $H^{\mathrm{MF}}$ 加均匀磁场：

\begin{equation*}
H ^ {\mathrm{MF}} \rightarrow H ^ {\mathrm{MF}} - h \sum_{is = \pm \frac {1}{2}} s \, e ^ {i \vec {\pi} \mathbf {X} _ {i}} a _ {is} ^ {\dagger} a _ {is}.
\tag{18.92}
\end{equation*}

允许约束场有均匀与交错分量：

\begin{equation*}
\lambda_ {i} = \lambda + \exp \left( i \vec {\pi} \mathbf {x} _ {i} \right) \lambda_ {s}.
\tag{18.93}
\end{equation*}

证明平均场基态能量在 $\lambda_{s} = -\frac{1}{2}h$ 时极小。讨论 $\lambda_{s}$ 如何可解释为 $xy$ 平面内所有自旋以角频率 $\omega = -\frac{1}{2}h$ 的均匀进动。

\item 利用上题结果，把 (18.63) 的均匀磁化率 $\chi_{0}^{\mathrm{MF}}$ 导出为自由能对 $h_{0}$ 的二阶导数。注意取热力学极限之前必须保持温度有限。为什么？

\end{exercises}

\begin{chapbib}
\item Schwinger 玻色子平均场理论引入于
  D. P. Arovas and A. Auerbach, Phys. Rev. B \textbf{38}, 316 (1988)。
\item 正方格子反铁磁体 SBMFT 动力学自旋关联见
  A. Auerbach and D. P. Arovas, Phys. Rev. Lett. \textbf{61}, 617 (1988)。
\item 一维经典 Heisenberg 铁磁体关联长度的计算：
  M. Fisher, Am. J. Phys. \textbf{32}, 343 (1964)。
\item 椭圆函数定义于
  M. Abramowitz and I. A. Stegun, \textit{Handbook of Mathematical Functions} (Dover, 1965), 第十七章。
\item 修正自旋波理论（含 SBMFT 方程的解析解）见
  M. Takahashi, Prog. Theor. Phys. Suppl. \textbf{87}, 233 (1986)；
  M. Takahashi, Phys. Rev. B \textbf{36}, 3791 (1987)；
  M. Takahashi, Phys. Rev. B \textbf{40}, 2494 (1989)。
\item 长程磁序与 Schwinger 玻色子凝聚的联系解释于
  J. E. Hirsch and S. Tang, Phys. Rev. B \textbf{39}, 2850 (1989)；
  M. Raykin and A. Auerbach, Phys. Rev. Lett. \textbf{70}, 3808 (1993)。
\item 大族 $SU(N)$ Heisenberg 模型的连续映射及拓扑 Berry 相位项的恢复由
  N. Read and S. Sachdev, Nucl. Phys. B \textbf{316}, 609 (1989)；
  N. Read and S. Sachdev, Phys. Rev. Lett. \textbf{62}, 1694 (1989)
  完成。
\end{chapbib}
